% mode: ini-etex
% Aimed at the effective tail, find_effective_tail and fetch_effective_tail: lists that end in the end_M_code node of a display
% line in TeXXeT mode, so that \lastbox, \unskip, \unkern and \unpenalty look at the node before it, and when only the
% begin_M/end_M pair is left, removal of the pair (fm=true).  And box numbers past 255: \box, \copy, \unhbox, \lastbox into them.
\input prelude
\output={\setbox0\box255 \deadcycles=0 }
\def\reset{\noindent\par}
\errorcontextlines=0
\setbox300\hbox{x}\showbox300
\setbox301\copy300 \showbox300 \showbox301
\setbox301\box300 \showbox300 \showbox301
\setbox302\copy400 \showbox302 \setbox302\box400 \showbox302
{\setbox303\box301 \showbox301 \showbox303 }\showbox301 \showbox303
\setbox0\vbox{\hbox{}\setbox304\lastbox\showbox304 }
% A display line with \predisplaydirection, set inside the display because it is computed at its start: app_display packages
% the line between math nodes.
\TeXXeTstate=1 \hsize=100pt
% Makes box 3 from the display's line: #1 is the direction, #2 the formula.
\def\display#1#2{\setbox1\vbox{\noindent\hbox{a}$$\predisplaydirection=#1 #2$$\par\unskip\unpenalty
  \setbox2\lastbox\global\setbox3\box2 }\showbox3 }
\display1{\hbox{b}}
% The list ends with end_M: each of the deletions and \lastbox looks through it.
\setbox5\hbox{\unhcopy3 \showlists\unskip\showlists\unskip\showlists}\showbox5
\setbox5\hbox{\unhcopy3 \unkern\showlists\unpenalty\showlists}\showbox5
\setbox5\hbox{\unhcopy3 \setbox6\lastbox\showbox6 \showlists}\showbox5
\setbox5\hbox{\unhcopy3 \unkern\setbox6\lastbox\showbox6 \showlists}\showbox5
\setbox5\hbox{\unhcopy3 \unkern\setbox6\lastbox\showbox6 \showlists\unkern\showlists\unkern\showlists}\showbox5
% With both kerns gone and the box gone, the walk finds begin_M in front of tx.
\setbox5\hbox{\unhcopy3 \unkern\setbox6\lastbox\unkern\showlists\unkern\showlists\setbox6\lastbox\showbox6 \showlists}\showbox5
% In the middle of a longer list: something before the display line's nodes.
\setbox5\hbox{x\hskip1pt\unhcopy3 \unkern\setbox6\lastbox\unkern\showlists\unkern\showlists}\showbox5
\display{-1}{\hbox{b}}
\setbox5\hbox{\unhcopy3 \showlists\unkern\setbox6\lastbox\showbox6 \showlists\unkern\showlists\unkern\showlists}\showbox5
\display1{}
\setbox5\hbox{\unhcopy3 \showlists\unkern\showlists\unkern\showlists\unkern\showlists}\showbox5
% With \leftskip and \rightskip the kerns are glue instead, cancelled by others.
\leftskip=3pt \rightskip=4pt
\display1{\hbox{b}}
\setbox5\hbox{\unhcopy3 \showlists\unskip\showlists\unkern\showlists\setbox6\lastbox\showbox6 \showlists\unskip\showlists\unskip\showlists}\showbox5
\display{-1}{\hbox{b}}
\setbox5\hbox{\unhcopy3 \showlists\unskip\unkern\setbox6\lastbox\showbox6 \showlists\unkern\unskip\showlists}\showbox5
\leftskip=0pt \rightskip=0pt
% With an equation number: a list of more than the box.
\display1{\hbox{b}\eqno\hbox{c}}
\setbox5\hbox{\unhcopy3 \showlists\unkern\setbox6\lastbox\showbox6 \showlists}\showbox5
% A plain \endL at the tail is not end_M_code: the tail stays what it is.
\setbox5\hbox{\beginL\hbox{}\endL\setbox6\lastbox\showbox6 \showlists}\showbox5
\setbox5\hbox{\beginL\hskip1pt\endL\unskip\showlists}\showbox5
\TeXXeTstate=0
\showlists
\end
